\section{Conclusion}
\label{sec:conclusion}

This paper presented UCRA, a closed-loop framework that converts
calibrated probabilistic demand forecasts into auditable capacity
reservations for 5G network slicing, together with a verified
experimental methodology for evaluating its self-evolving component
honestly. The framework's decision core is a single transform,
$\Rt = \qtau + \kappa\,\spread\,(1 + w\,\rhot)$, which consumes a
deep quantile forecast, prices forecast dispersion with one operator
dial $\kappa$, and modulates the buffer with an online risk indicator
--- wrapped in EWMA smoothing with release hysteresis, and backed by
a drift-triggered, reservoir-replay fine-tuning stage operating under
a formally causal information-flow contract.

The verified evidence supports four conclusions. First, on three
weeks of live RAN performance counters, the closed loop holds 0.0\%
reservation violations at 72.1\% utilization with 10.7\% waste ---
against 30.5\% waste for static peak sizing, 48.0\% violations for
median-forecast sizing, and 16.3\% violations for an oracle quantile
that cheats with the test window's own future --- and the result is
stable across seeds (0.0\% violations in all; utilization
$67.6\% \pm 3.4$). Second, the kappa dial traces a clean empirical
risk--utilization frontier whose zero-violation knee lies between
$\kappa = 0$ and $0.25$ on this data, turning SLA policy into a
measured positioning decision rather than a guess. Third, under a
$+35\%$ demand surge, the causal self-evolution stage halves
post-surge violations ($34.3\% \pm 5.9 \rightarrow 17.2\% \pm 4.9$)
while firing zero spurious updates on clean data across all seeds
--- the symmetric silence-when-unneeded, act-when-needed behavior
that trigger-based adaptivity should be held to. Fourth, on the
independent CTTC slicing workload the same transform reduces slice
violations to 1.2--2.7\% across URLLC, eMBB, and mMTC where the
operator's default sizing violates 5.8--43.7\% --- with the honest
finding that on that dataset the buffer term, not the risk
modulation, does the work.

Beyond the numbers, the methodological contribution may outlast them.
Our own pre-submission audit found that the self-evolution evaluation
leaked future labels through the replay path, refreshed served
forecasts retroactively, and let a CTTC policy see the outcome it
protected against; the final pipeline fixes all three, proves the
fix sufficient, pins it with regression tests, and quantifies how
the corrected protocol changes the headline results. We release the
full pipeline --- ingestion, training, closed loop, causality tests,
and the three-seed verification driver --- so that the community can
re-run, re-audit, and extend it. Future work proceeds along four
tracks. \emph{Evidence breadth}: multi-segment evaluation on the
dataset's remaining two networks, longer windows covering seasonal
transitions, and natural-drift (rather than injected) adaptation
cases. \emph{Metric depth}: severity-graded SLA metrics (violation
duration, depth, and affected-user counts) replacing the binary
per-slot flag, and slice-level time-series evaluation of Stage~3's
per-slice risk weighting. \emph{Controller comparison}: benchmarking
UCRA's auditable transform against reinforcement-learning slice
controllers under the same causal evaluation contract --- a
comparison we expect to be informative precisely because the
contract neutralizes the information advantages that flatter learned
policies. \emph{Mechanism}: replacing the fixed replay delay with a
learned-availability schedule, and studying trigger co-design ---
threshold, cadence, and fine-tune intensity --- as a joint
risk-allocation problem.
